\section*{NeurIPS Paper Checklist}

\begin{enumerate}

\item {\bf Claims}
    \item[] Question: Do the main claims made in the abstract and introduction accurately reflect the paper's contributions and scope?
    \item[] Answer: \answerYes{}
    \item[] Justification: The abstract states five specific claims (TF-IDF outperforms neural methods on BackdoorLLM, signal is in response content, custom code backdoors evade all methods, orthogonality of behavioral and structural directions, causal verification at 90\% precision). Each is directly supported by experimental results in Sections 4.1--4.6. The scope is clearly bounded: we evaluate on BackdoorLLM and custom code backdoors, not on all possible attack types. Limitations are explicitly discussed in Section 5.

\item {\bf Limitations}
    \item[] Question: Does the paper discuss the limitations of the work performed by the authors?
    \item[] Answer: \answerYes{}
    \item[] Justification: Section 5 (Discussion) includes a dedicated Limitations paragraph covering: small corpus size (800 documents), degradation with heterogeneous data sources (Section 4.6, scale test), limited cross-architecture evaluation (one additional family), 4-bit quantization effects on 13B/70B results, single-seed causal and adversarial experiments, untested pre-training poisoning, and the pending need for more diverse subtle attack designs.

\item {\bf Theory assumptions and proofs}
    \item[] Question: For each theoretical result, does the paper provide the full set of assumptions and a complete (and correct) proof?
    \item[] Answer: \answerNA{}
    \item[] Justification: This is an empirical paper. We do not present theoretical results, theorems, or proofs. All claims are supported by experimental evidence.

    \item {\bf Experimental result reproducibility}
    \item[] Question: Does the paper fully disclose all the information needed to reproduce the main experimental results of the paper to the extent that it affects the main claims and/or conclusions of the paper (regardless of whether the code and data are provided or not)?
    \item[] Answer: \answerYes{}
    \item[] Justification: Section 3 specifies all models (with HuggingFace IDs), datasets (BackdoorLLM with document counts), detection methods (10 methods with exact specifications), hyperparameters (LoRA rank 16, alpha 32, 5 epochs, lr $2 \times 10^{-4}$, seeds 42/123/456), quantization settings (4-bit NF4 via bitsandbytes), evaluation metric (AUROC, formally defined), and infrastructure (H100 80GB, PyTorch 2.8, transformers 4.57, peft 0.13). The custom code backdoors are described with trigger phrases, target behaviors, and training recipes. TF-IDF uses scikit-learn with max\_features=5000.

\item {\bf Open access to data and code}
    \item[] Question: Does the paper provide open access to the data and code, with sufficient instructions to faithfully reproduce the main experimental results, as described in supplemental material?
    \item[] Answer: \answerYes{}
    \item[] Justification: All experiment scripts, result CSVs, and V matrices are available in the supplementary material. BackdoorLLM models and datasets are publicly available on HuggingFace (BackdoorLLM organization). CodeAlpaca-20K is publicly available. Code includes modular scripts for each experiment with a master orchestrator. A README documents the directory structure and how to reproduce each result.

\item {\bf Experimental setting/details}
    \item[] Question: Does the paper specify all the training and test details (e.g., data splits, hyperparameters, how they were chosen, type of optimizer) necessary to understand the results?
    \item[] Answer: \answerYes{}
    \item[] Justification: Section 3 specifies: data splits (400 poison + 400 clean per attack, subsampled at 5 poison rates with 3 seeds), optimizer (AdamW with lr $2 \times 10^{-4}$), LoRA configuration (rank 16, alpha 32, dropout 0.05, target modules q\_proj/v\_proj), training epochs (5), batch size (2), sequence length (256 for activation extraction, 512 for training), TF-IDF settings (max\_features=5000, English stop words), sentence transformer model (all-MiniLM-L6-v2), and LinearDCT settings (64 factors, 8 layers, 300 training documents, 10 VJP iterations).

\item {\bf Experiment statistical significance}
    \item[] Question: Does the paper report error bars suitably and correctly defined or other appropriate information about the statistical significance of the experiments?
    \item[] Answer: \answerYes{}
    \item[] Justification: The poison rate sweep (Appendix A) reports mean $\pm$ standard deviation over 3 random seeds for all 5 attacks. The detection ladder reports mean AUROC over 5 attacks with the standard deviation noted in the text (std = 0.002 for TF-IDF, std = 0.002 for L2-blind). The causal verification and adversarial experiments are single-seed, which is explicitly noted in the text and figure captions. Cross-attack variance is the primary source of variability and is reported throughout.

\item {\bf Experiments compute resources}
    \item[] Question: For each experiment, does the paper provide sufficient information on the computer resources (type of compute workers, memory, time of execution) needed to reproduce the experiments?
    \item[] Answer: \answerYes{}
    \item[] Justification: Section 3 specifies NVIDIA H100 80GB GPUs. Section 4.6 provides per-method runtime: TF-IDF 0.1 ms/doc, sentence embeddings 0.8 ms/doc, L2-blind 8.4 ms/doc, perplexity 7.8 ms/doc. The full experimental suite (1,200+ conditions across BackdoorLLM plus custom backdoors) ran on two 8-GPU H100 pods over approximately 8 hours. The total compute for the project, including preliminary experiments, was approximately 200 GPU-hours on H100.

\item {\bf Code of ethics}
    \item[] Question: Does the research conducted in the paper conform, in every respect, with the NeurIPS Code of Ethics \url{https://neurips.cc/public/EthicsGuidelines}?
    \item[] Answer: \answerYes{}
    \item[] Justification: This work studies backdoor detection, which is a defensive security capability. We use publicly available benchmark models (BackdoorLLM) that were released specifically for security research. Our custom code backdoors are trained on standard code instruction datasets and demonstrate insecure coding patterns (not novel attack vectors). We do not release backdoored models that could be misused. The research aims to improve the security of LLM training pipelines.

\item {\bf Broader impacts}
    \item[] Question: Does the paper discuss both potential positive societal impacts and negative societal impacts of the work performed?
    \item[] Answer: \answerYes{}
    \item[] Justification: Section 5 discusses positive impacts (enabling pre-training data screening, practical recommendations for practitioners) and negative implications (our demonstration that subtle backdoors evade all detection could inform attackers, though the attack techniques we use are already published). The paper's primary message is cautionary: current detection creates a false sense of security because benchmarks are too easy.

\item {\bf Safeguards}
    \item[] Question: Does the paper describe safeguards that have been put in place for responsible release of data or models that have a high risk for misuse (e.g., pre-trained language models, image generators, or scraped datasets)?
    \item[] Answer: \answerNA{}
    \item[] Justification: We do not release pre-trained backdoored models. The custom code backdoors are trained with standard LoRA on publicly available data and demonstrate known vulnerability patterns (insecure template rendering, missing input validation), not novel attack capabilities. The BackdoorLLM models we use are already publicly available on HuggingFace.

\item {\bf Licenses for existing assets}
    \item[] Question: Are the creators or original owners of assets (e.g., code, data, models), used in the paper, properly credited and are the license and terms of use explicitly mentioned and properly respected?
    \item[] Answer: \answerYes{}
    \item[] Justification: BackdoorLLM is cited and its models are accessed via HuggingFace under their published terms. LLaMA-2 models are accessed via NousResearch community mirrors. CodeAlpaca-20K is cited. Mistral-7B and Gemma-2-2B are accessed via their official HuggingFace repositories. LinearDCT, CAA, and SPECTRE methods are cited with their original papers. scikit-learn, PyTorch, transformers, peft, and bitsandbytes are standard open-source libraries.

\item {\bf New assets}
    \item[] Question: Are new assets introduced in the paper well documented and is the documentation provided alongside the assets?
    \item[] Answer: \answerYes{}
    \item[] Justification: We introduce custom code backdoor datasets (two variants with documented triggers, targets, and training recipes) and a detection evaluation framework (10 methods, documented in code). Both are provided in the supplementary material with a README explaining the directory structure, how to run each experiment, and how to interpret the output CSVs.

\item {\bf Crowdsourcing and research with human subjects}
    \item[] Question: For crowdsourcing experiments and research with human subjects, does the paper include the full text of instructions given to participants and screenshots, if applicable, as well as details about compensation (if any)?
    \item[] Answer: \answerNA{}
    \item[] Justification: This paper does not involve crowdsourcing or research with human subjects.

\item {\bf Institutional review board (IRB) approvals or equivalent for research with human subjects}
    \item[] Question: Does the paper describe potential risks incurred by study participants, whether such risks were disclosed to the subjects, and whether Institutional Review Board (IRB) approvals (or an equivalent approval/review based on the requirements of your country or institution) were obtained?
    \item[] Answer: \answerNA{}
    \item[] Justification: This paper does not involve human subjects research.

\item {\bf Declaration of LLM usage}
    \item[] Question: Does the paper describe the usage of LLMs if it is an important, original, or non-standard component of the core methods in this research? Note that if the LLM is used only for writing, editing, or formatting purposes and does \emph{not} impact the core methodology, scientific rigor, or originality of the research, declaration is not required.
    \item[] Answer: \answerYes{}
    \item[] Justification: LLMs (LLaMA-2, Mistral, Gemma) are central to the research as both the subject of backdoor attacks and the feature extractors used for detection. Their usage is fully described in Section 3 with model IDs, quantization settings, and layer specifications. Claude Code was used as a research assistant for experiment orchestration and code development.

\end{enumerate}
